\chapter{Introduction}
\label{chap:introduction}

\section{Problem Statement}

This section will outline the problem this dissertation seeks to address, 
and the context surrounding it.

According to \cite{hea2024key}, over 48,000 Irish secondary school students progressed to college after graduation in 2024.
The individuals who have considered progressing to college have all been posed with the question:
\emph{What college course should I study?} 
This is a daunting question for many reasons.
The majority of secondary school students are faced with making this decision 
at 16-19 years old, an age with limited life experience and knowledge about the world.
This decision also comes with many personally confronting questions, which may include:
\emph{``What do I want to do with the rest of my life?"} and \emph{``What do I love doing?"}, among many others.
Moreover, an individual may choose a college course they will later regret.
For instance, in \cite{hea2018graduate}, about 10,000 (31\%) of recent Irish college graduates reported they were not likely to study their course again, 
adding more weight to this crucial life decision.

Furthermore,
individuals are shown to be influenced by various socio-demographic factors when choosing their college course.
For instance in \cite{mullen2011degrees, mullen2003goes, leighton2023rich}, 
college course choice is largely influenced by the parent's educational background
of the student.
Additionally, \cite{slattery2023navigating, iwish2025survey} demonstrate the downsides of 
gender representation in certain college courses, acting as barriers to college students of 
under-represented genders.
Lastly, \cite{racial2015peter, dickson2010race, xie2003social} highlight the influence of 
race and ethnicity on an individual's college course choice.
The above points illustrate the fact that
 this decision is not made in a vacuum,
adding complexity to a decision that is already challenging.

To assist students in this difficult decision, it would be appropriate to introduce the role of 
a Guidance Counsellor (GC).
GCs are trained professionals in Irish secondary schools, primarily known for assisting students in their career prospects after graduation.
Most notably, the role of a GC extends well beyond vocational support.
According to \cite{igc2016core}, GCs are also trained to provide 
educational, personal and social support.
In this context, it is worth introducing the Institute of Guidance Counsellors (IGC)\footnote{\url{https://igc.ie/}}, the professional body for GCs in Ireland.
In a report by the \cite{igc2022survey}, mental health issues were more commonly encountered by GCs compared to
career decision-making issues\footnote{In \cite{igc2024survey}, career decision-making issues did surpass mental health issues as the most commonly encountered issue for GCs; however, both are pressing issues.}, marking the first time that career decision-making issues were not the most commonly encountered issue by GCs since 2019.
In the IGC's 2024-2027 strategic plan outlined in \cite{igc2024plan}, 
they signalled a shift in priorities for GCs across Irish secondary schools: vocational support may not be the most important aspect of a GCs role going forward.

As is evident, choosing a college course is a difficult, multi-faceted decision, and the decision is made more challenging 
by the shifting priorities of Irish GCs.
This leads us to the motivation behind this research.

\section{Research Motivation}
\label{sec:motivation}

In light of the identified problems, 
we can shift to a more technological focus. A Recommender System (RS) is a technology that recommends items to a user.
Well-known examples of RSs include Netflix\footnote{\url{https://www.netflix.com/}} recommending films to their customers, or Amazon\footnote{\url{https://www.amazon.com/}} suggesting products to their customers.
One of the earliest known RSs was \emph{Tapestry} in \cite{goldberg1992using},
a system that made tailored mailing list recommendations to its users. 
Around this time, the field of RSs received considerable interest as e-commerce websites could increase
their revenue if they suggested more personalised products to their users.
As discussed in \cite{schafer1999recommender}, Amazon and Netflix are the most notable adopters of RSs
in order to increase company revenue.

As previously explained, an RS recommends items to a user.
However, we can re-conceptualise an RS in the context of college courses:
college courses can take the place of items, and a student can take the place of a user.
Given the established difficulty and influences of an individual's college course choice, the shifting priorities of GCs towards student's mental health as opposed to career-making decisions,
and a technology that could be appropriately applied to this domain,
an opportunity presents itself: a Recommender System for Irish college courses.

A technology, such as an RS, could assist an individual in navigating a difficult decision-making process.
For instance, there are over 45 higher education institutes in Ireland offering a total of over 1700 courses\footnote{\label{}Higher education colleges and courses can be found here: \url{https://careersportal.ie/}},
representing a huge set of choice to the individual.
An advantage of a personalised RS is to provide a user with a much smaller list of relevant recommendations.
\cite{long2025choice} demonstrated that smaller recommendation sets are often equally as effective as larger recommendations sets,
suggesting that a much narrower set of personalised college course recommendations could be effective in this space.

For the 2026/2027 academic year, 52\footnote{See \url{https://careersportal.ie/courses/coursefinder?types_in=2} for new courses} new higher education courses 
were introduced in Ireland.
Moreover, as discussed in section \ref{sec:lc-points}, the requirements for Irish higher education courses 
change every year as a reflection of demand and changing priorities.
This introduces administrative strain on GCs, making it difficult to have up-to-date awareness on the newest courses 
and current course entry requirements.
A technology, such as an RS, can be automated to maintain the latest college course information
and utilise this information in recommending college courses that a human GC may not even be aware of, 
posing a potential advantage of an RS in this domain.

As we have argued, an RS is a technology that could be appropriately applied to the domain of Irish college courses,
potentially alleviating the difficulty of the college course decision-making process.
With the motivation behind our research described, we can now define our research question.

\section{Research Question}
\label{sec:research-question}

In light of the motivation behind this research, the research question for this dissertation can be formulated:
\emph{``To what extent can a Recommender System assist an Irish student in exploring college courses they would consider studying?"}

In this research question, the meaning of `assist' is intentionally broad, 
and we can further understand the meaning of `assist' by exploring the three qualities of an effective RS, as outlined in \cite{aggarwal2016recommender}, a prominent and popular textbook on RSs:

According to Aggarwal, the recommendations should be \textbf{relevant} to the user's preferences. 
In this context, an effective RS would recommend college courses the user would consider studying based off their gathered preferences.
For example, it would generally not make sense to recommend healthcare courses to a student that has demonstrated no interest in healthcare.
In context to our research question, an RS with relevant recommendations would assist a student by recommending college courses they would consider studying.

Aggarwal also argues that the recommendations should be \textbf{diverse} and avoid being uniform in a particular category.
For example, an effective RS would not suggest 20 different computer science courses to a user, even if the user has demonstrated 
a strong interest in computers. If we revisit the research question, we desire an RS that assists a student's college course exploration,
and by suggesting college courses across a diverse range of the user's preferences,
the RS would assist the user in discovering college courses they would consider studying.

A user may receive recommendations that are relevant to their preferences and diverse across their interests,
however the user may already have an awareness of these recommendations.
For example, a student may receive recommendations in biology and engineering, but they may have already considered studying these college courses previously.
As a result, the user does not actually gain any new information after interacting with the RS.
Aggarwal claims that the quality that separates a good RS from a great one is \textbf{serendipity:}
these are recommendations that are genuinely unexpected but incredibly relevant.
The RS demonstrates a deep knowledge of the user, providing the user with relevant recommendations that the user was genuinely unaware of.
Continuing our example, a serendipitous recommendation to this user might be biomedical engineering, merging both of the user's interests in an unexpected but useful way,
and the user may now be considering studying the college course after interacting with the RS.
This would be incredibly effective in assisting the user with their college course exploration, as outlined in the research question.

Serendipity is the gold standard for RSs.
However, it is very challenging to consistently produce serendipitous recommendations,
because by their nature the recommendations are `high-risk, high-reward'.
RSs are often faced with the dilemma between producing `safe' recommendations that are likely relevant but uninteresting,
or more `adventurous' recommendations that could be serendipitous or unsatisfactory.
This is a well-known tradeoff in the field of RSs, and it is the price paid for attempting to achieve the gold standard in RSs: serendipity.

In summary, according to \cite{aggarwal2016recommender}, an effective RS provides recommendations that are:

\begin{itemize}
    \item Relevant
    \item Diverse
    \item Serendipitous
\end{itemize}

It is worth mentioning that the above three qualities will be referenced frequently throughout this dissertation, 
and thus we recommend that the reader grows familiar with these concepts.
In summary, the research question for this dissertation is: \emph{``To what extent can a Recommender System assist an Irish student in exploring college courses they would consider studying?"},
and we hypothesise that an RS that provides relevant, diverse, and serendipitous recommendations would assist an Irish student in exploring college courses they would consider studying.
Therefore, this dissertation's ability to answer this research question will be measured by its performance on the relevance, diversity and serendipity metrics.

\section{Research Objectives}

With our research question formulated, we can define our research objectives below:
\begin{enumerate}
    \item \label{res-obj-answer} Provide an answer to the question: \emph{``What college course should people study?"}.
    \item \label{res-obj-information} Identify the relevant information an RS should question a user to recommend suitable college courses.
    \item \label{res-obj-baseline} Design an objective, easily reproducible baseline RS appropriate to the Irish college course context.
    \item \label{res-obj-trust} In comparison to a baseline college course RS, design an RS with statistically significant improvements for the user-perceived trust metric.
    \item \label{res-obj-insights} Analyse the qualitative insights and findings from a conducted experiment.
\end{enumerate}

\section{Overview}

In this section, an overview of this dissertation's chapters is provided.

\begin{itemize}
    \item Chapter \ref{chap:background-reading} - Background Reading. This chapter will describe the relevant background knowledge to understand the Irish college admission system.
    \item Chapter \ref{chap:literature-review} - Literature Review. This chapter will review the literature surrounding college course choice, and closely-related works.
    \item Chapter \ref{chap:design} - Design. This chapter will describe the high-level design of the college course RS, with a justification of different design decisions.
    \item Chapter \ref{chap:implementation} - Implementation. This chapter will explore the implementation details behind the design ideas exposed in chapter \ref{chap:design}.
    \item Chapter \ref{chap:evaluation} - Evaluation. This chapter will explain the conducted experiment methodology, and evaluate the RS's performance on various metrics.
    \item Chapter \ref{chap:conclusions} - Conclusions \& Future Work. This chapter will summarise the dissertation, and explore potential areas for future development.
\end{itemize}
